\section{The \texttt{traceq} Library}
\label{sec:mmc-traceq}

%4.4 The \texttt{traceq} Library
%• 4.4.1 Motivation & Design Goals
%• Why a new library was created (e.g., bridging the gap between external checks and Python-internal traces, providing an easy-to-use API).
%• Outline key features or objectives (lightweight overhead, integration with HPC workflows, easy Docker usage, etc.).
%• 4.4.2 Implementation Overview
%• Explain the technical backbone of traceq:
%• Internally uses tracemalloc or other measurement modules.
%• Possibly collects Docker stats or OS-level info.
%• Summarizes memory usage in a single place (peak usage, usage over time).
%• Discuss how it interacts with the Python GC or hooking into process metrics.
%• 4.4.3 Usage Examples
%• Show a minimal code snippet: how to import and start traceq, how to record memory usage, how to interpret the results.
%• Potential example: traceq used in a typical HPC pipeline or a short script demonstrating real-time memory logging.
%• 4.4.4 Validation & Comparisons
%• Summarize any testing or benchmark used to verify traceq results against known references (e.g., Docker stats, psutil, or the memory constraints method).
%• If relevant, mention performance overhead: e.g., “Our experiments showed overhead <5% in CPU and memory usage.”